\chapter{Completion of the proof of Theorem~{\ref{MAIN}}}

We have established the requisite non-collapsing result assuming the
existence of strong canonical neighborhoods. In order to complete
the proof of Theorem~\ref{MAIN} it remains for us to show the
existence of strong canonical neighborhoods. This is the result of
the next section.


\section{Proof of the strong canonical neighborhood assumption}\label{sectcannbhd}


 \begin{proposition}[The strong canonical neighborhood assumption extends to the next scale r_{i+1}]
\label{prop:canonical-neighborhood-extends}
\uses{prop:kappa-limit-noncollapsing,thm:main-existence-ricci-flow-surgery}\label{extend}
Suppose that for some $i\ge 0$ we have surgery parameter sequences
$\delta_0\ge \delta_1\ge\cdots\ge \delta_i>0$, $\epsilon=r_0\ge
r_1\ge \cdots\ge r_i>0$ and $\kappa_0\ge \kappa_1\ge \cdots\ge
\kappa_i>0$. For any $r_{i+1}\le r_i$, let $\delta(r_{i+1})>0$ be
the constant in Proposition~\ref{KAPPALIMIT} associated to these
three sequences and to $r_{i+1}$.
 Then there are positive
constants $r_{i+1}\le r_i$ and $\delta_{i+1}\le\delta(r_{i+1})$ such that the
following holds. Suppose that $({\mathcal M},G)$ is a Ricci flow with surgery
defined for $0\le t<T$ for some $T\in (T_i,T_{i+1}]$ with surgery control
parameter $\overline\delta(t)$. Suppose that the restriction of this Ricci flow
with surgery to ${\bf t}^{-1}([0,T_i))$ satisfies Assumptions (1) -- (7) and
also the five properties given in the hypothesis of Theorem~\ref{MAIN} with
respect to the given sequences. Suppose also that $\overline\delta(t)\le
\delta_{i+1}$ for all $t\in [T_{i-1},T]$. Then $({\mathcal M},G)$  satisfies
the strong $(C,\epsilon)$-canonical neighborhood assumption with parameter
$r_{i+1}$.
 \end{proposition}


\begin{proof}
\uses{prop:kappa-limit-noncollapsing,kaplimit,limitcannbhd,sliceconv,canonvary,narrows,controlnbhd}
Suppose that the result does not hold.
  Then we can take a sequence of $r_a\rightarrow 0$ as
 $a\rightarrow \infty$, all less than $r_i$, and for each $a$ a sequence
 $\delta_{a,b}\rightarrow 0$ as $b\rightarrow \infty$ with each
 $\delta_{a,b}\le \delta(r_a)$, where $\delta(r_a)\le \delta_i$ is the constant in Proposition~\ref{KAPPALIMIT}
 associated to the three sequences given in the statement of this proposition and $r_a$, such that for
 each $a,b$ there is a Ricci flow with surgery $({\mathcal
 M}_{(a,b)},G_{(a,b)})$ defined for $0\le t<T_{(a,b)}$ with $T_i<
 T_{(a,b)}\le T_{i+1}$ with control parameter $\overline \delta_{(a,b)}(t)$
 such the flow satisfies the hypothesis of the proposition with
 respect to these constants but fails to satisfy the conclusion.

\begin{lemma}[A limiting counterexample sequence of failure points converges to a limit failure point]
\label{lem:limit-canonical-neighborhood-fails-at-surgery-scale}\label{16.2}
For each $a$, and given $a$, for all $b$ sufficiently large there is $t_{(a,b)}\in
[T_i,T_{(a,b)})$ such that the restriction of $({\mathcal
M}_{(a,b)},G_{(a,b)})$ to ${\bf t}^{-1}\left([0,t_{(a,b)})\right)$
satisfies the strong $(C,\epsilon)$-canonical neighborhood
assumption with parameter $r_a$ and furthermore, there is
$x\in{\mathcal M}_{(a,b)}$ with ${\bf t}(x_{(a,b)})=t_{(a,b)}$ at
which the strong $(C,\epsilon)$-canonical neighborhood assumption
with parameter $r_a$ fails.
\end{lemma}

\begin{proof}
\uses{thm:local-surgery,prop:long-L-length-near-surgery-cap,lem:surgery-cap-ball-bounds,defncanonnbhd,defnepsround,Ccomponent,canonvary,neckcurv,A_0,Restim,stdsoln,bigell,neckglue,stdsolncannbhd}
Fix $a$.
 By supposition, for each $b$ there is a point $x\in {\mathcal M}_{(a,b)}$ at
which the strong $(C,\epsilon)$-canonical neighborhood assumption
fails for the parameter $r_a$. We call points at which this
condition fails {\em counterexample points}. Of course, since
$r_a\le r_i$ and since the restriction of $({\mathcal
M}_{(a,b)},g_{(a,b)})$ to ${\bf t}^{-1}([0,T_i))$ satisfies the
hypothesis of the proposition, we see that any counterexample point
$x$ has ${\bf t}(x)\ge T_i$. Take a sequence $x_n=x_{n,(a,b)}$   of
counterexample points with ${\bf t}(x_{n+1})\le {\bf t}(x_{n})$ for
all $n$ that minimizes ${\bf t}$ among all counterexample points in
the sense that for any $\xi>0$ and for any counterexample point
$x\in {\mathcal M}_{(a,b)}$ eventually ${\bf t}(x_{n})< {\bf
t}(x)+\xi$. Let $t'=t_{(a,b)}'=\mathrm{lim}_{n\rightarrow\infty}{\bf
t}(x_{n})$. Clearly, $t'\in [T_i,T_{(a,b)})$, and by construction
the restriction of $({\mathcal M}_{(a,b)},G_{(a,b)})$ to ${\bf
t}^{-1}([0,t'))$ satisfies the $(C,\epsilon)$-canonical neighborhood assumption
with parameter $r_a$. Since the surgery times are discrete, there is
$t''=t_{(a,b)}''$ with $t'<t''\le T_{(a,b)}$ and a diffeomorphism
$\psi=\psi_{(a,b)}\colon M_{ t'}\times[t',t'')\to {\bf
t}^{-1}([t',t''))$ compatible with time and the vector field. We
view $\psi^*G_{(a,b)}$ as a one-parameter family of metrics
$g(t)=g_{(a,b)}(t)$ on $M_{t'}$ for $t\in [t',t'')$. By passing to a
subsequence we can arrange that ${\bf t}(x_{n})\in [t',t'')$ for all
$n$. Thus, for each $n$ there are $y_n=y_{n,(a,b)}\in M_{t'}$ and
$t_{n}\in [t',t'')$ with $\psi(y_{n},t_{n})=x_{n}$. Since
$M_{t'}$ is a compact $3$-manifold, by passing to a further
subsequence we can assume that $y_{n}\rightarrow x_{(a,b)}\in
M_{t'}$. Of course, $t_{n}\rightarrow t'$ as $n\rightarrow\infty$ and
 $\mathrm{lim}_{n\rightarrow
\infty}x_{n}=x_{(a,b)}$ in ${\mathcal M}_{(a,b,)}$.

 We claim that, for
all $b$ sufficiently large, the strong $(C,\epsilon)$-canonical neighborhood assumption
with parameter $r_a$ fails at $x_{(a,b)}$. Notice that since $x_{(a,b)}$
is the limit of a sequence where the strong
$(C,\epsilon)$-neighborhood assumption fails, the points in the
sequence converging to $x_{(a,b)}$ have scalar curvature at least
$r_a^{-2}$. It follows that $R(x_{(a,b)})\ge r_a^{-2}$.  Suppose
that $x_{(a,b)}$ satisfies the strong $(C,\epsilon)$-canonical
neighborhood assumption with parameter $r_a$. This means that there
is a neighborhood $U=U_{(a,b)}$ of $x_{(a,b)}\in M_{t'}$
which is a strong $(C,\epsilon)$-canonical neighborhood of
$x_{(a,b)}$. According to Definition~\ref{defncanonnbhd} there are
four possibilities. The first two we consider are that $(U,g(t'))$
is an $\epsilon$-round component or a $C$-component. In either of
these cases, since the defining inequalities given in
Definition~\ref{defnepsround} and~\ref{Ccomponent} are strong
inequalities,  all metrics on $U$ sufficiently close to $g(t')$
in the $C^\infty$-topology the satisfy these same inequalities. But as $n$ tends to
$\infty$, the metrics $g(t_n)|_U$ converge in the $C^\infty$-topology
to $g(t')|_U$. Thus, in these two cases, for all $n$ sufficiently
large, the metrics $g(t_n)$ on $U$ are $(C,\epsilon)$-canonical
neighborhood metrics of the same type as $g(t_{(a,b)}')|_U$. Hence, in
either of these cases, for all $n$ sufficiently large $x_{n,(a,b)}$
has a strong $(C,\epsilon)$-canonical neighborhood of the same type
as $x_{(a,b)}$, contrary to our assumption about the sequence
$x_{n,(a,b)}$.



Now suppose that there is
 a $(C,\epsilon)$-cap whose core contains $x_{(a,b)}$.
This is to say that $(U,g(t'))$ is a $(C,\epsilon)$-cap whose core
contains $x_{(a,b)}$. By Proposition~\ref{canonvary}, for all $n$
sufficiently large, $(U,g(t_{n}))$ is also a $(C,\epsilon)$-cap with
the same core. This core contains $y_{n}$ for all $n$ sufficiently
large, showing that $x_{n}$ is contained in the core of a
$(C,\epsilon)$-cap  for all $n$ sufficiently large.


Now let us consider the remaining case when $x_{(a,b)}$ is the
center of a strong $\epsilon$-neck. In this case we have an
embedding $\psi_{U_{(a,b)}}\colon
U_{(a,b)}\times({t_{(a,b)}'}-R^{-1}(x_{(a,b)}),t_{(a,b)}']\to
{\mathcal M}_{(a,b)}$ compatible with time and the vector field and
a diffeomorphism $f_{(a,b)}\colon S^2\times
(-\epsilon^{-1},\epsilon^{-1})\to U_{(a,b)}$ so that $
(f_{(a,b)}\times\Id)^*\psi_{U_{(a,b)}}^*(R({x_{(a,b)}})G_{(a,b)})$ is
$\epsilon$-close in the $C^{[1/\epsilon]}$-topology to the
evolving product metric $h_0(t)\times ds^2,\ -1<t\le 0$, where $h_0(t)$
is a round metric of scalar curvature $1/(1-t)$ on $S^2$ and $ds^2$
is the Euclidean metric on the interval. Here, there are two
subcases to consider.
\begin{enumerate}
\item[(i)] $\psi_{U_{(a,b)}}$ extends backward past $
t_{(a,b)}'-R^{-1}(x_{(a,b)})$.
\item[(ii)] There is a flow line through a point $y_{(a,b)}\in U_{(a,b)}$ that is defined
on the interval $[t_{(a,b)}'-R^{-1}(x_{(a,b)}),t_{(a,b)}']$ but with
the value of the flow line at $t_{(a,b)}'-R^{-1}(x_{(a,b)})$ an
exposed point.
\end{enumerate}

Let us consider the first subcase. The embedding $\psi_{U_{(a,b)}}$
extends forward in time because of the diffeomorphism
$\psi_{(a,b)}\colon M_{t_{(a,b)}'}\times [t_{(a,b)}',t_{(a,b)}'')\to
{\mathcal M}_{(a,b)}$ and, by assumption, $\psi_{U_{(a,b)}}$ extends
backward in time some amount. Thus, for all $n$ sufficiently large,
we can use these extensions of $\psi_{U_{(a,b)}}$ to define an
embedding $\psi_{n,(a,b)}\colon U_{(a,b)}\times ({\bf
t}(x_{n,(a,b)})-R^{-1}(x_{n,(a,b)}),{\bf t}(x_{n,(a,b)})]\to
{\mathcal M}_{(a,b)}$ compatible with time and the vector field.
Furthermore, since the $\psi_{n,(a,b)}$ converge in the
$C^\infty$-topology as $n$ tends to infinity to $\psi_{U_{(a,b)}}$,
the Riemannian metrics $(f_{(a,b)}\times\Id)^*\psi_{n,(a,b)}^*(R(x_{n,(a,b)})G_{a,b})$ converge in the
$C^\infty$-topology to the pullback  $(f_{(a,b)}\times\Id)^*\psi_{U_{(a,b)}}^*(R(x_{(a,b)})G_{a,b})$. Clearly then, for
fixed $(a,b)$ and for all $n$ sufficiently large the pullbacks of
the rescalings of these metrics by $R(x_{n,(a,b)})$ are within
$\epsilon$ in the $C^{[1/\epsilon]}$-topology of the standard
evolving flow $h_0(t)\times ds^2, -1<t\le 0$, on the product of
$S^2$ with the interval. Under these identifications the points
$x_{n,(a,b)}$ correspond to points $(p_{n,(a,b)},s_{n,(a,b)})\in
S^2\times (-\epsilon^{-1},\epsilon^{-1})$ where $\mathrm{lim}_{n\rightarrow\infty}s_{n,(a,b)}=0$. The last thing we  do is
to choose diffeomorphisms $\varphi_{n,(a,b)}\colon
(-\epsilon^{-1},\epsilon^{-1})\to (-\epsilon^{-1},\epsilon^{-1})$
that are the identity near both ends, such that $\varphi_{n,(a,b)}$
carries $0$ to $s_{n,(a,b)}$ and such that the $\varphi_{n,(a,b)}$
converge to the identity in the $C^\infty$-topology for fixed
$(a,b)$ as $n$ tends to infinity. Then, for all $n$ sufficiently
large, the composition
$$S^2\times (-\epsilon^{-1},\epsilon^{-1})\stackrel{\Id\times
\varphi_{n,(a,b)}}{\longrightarrow}S^2\times
(-\epsilon^{-1},\epsilon^{-1})\stackrel{f_{(a,b)}}{\longrightarrow}
U\stackrel{\psi_{n,(a,b)}}{\longrightarrow} {\mathcal M}_{(a,b)}$$
is a strong $\epsilon$-neck centered at $x_{n,(a,b)}$. This shows that for any
$b$ for which the first subcase holds, for all $n$ sufficiently large, there is
a strong $\epsilon$-neck centered at $x_{n,(a,b)}$.

Now suppose that the second subcase holds for all $b$. Here, unlike
all previous cases, we shall have to let $b$ vary and we shall prove
the result only for $b$ sufficiently large.
We shall show that for all $b$ sufficiently large, $x_{(a,b)}$ is contained
in the core of a $(C,\epsilon)$-cap. This will establish the result by contradiction,
for as we showed in the previous case, if $x_{(a,b)}$ is contained in the core of
a $(C,\epsilon)$-cap, then the same is true for the $x_n$ for all $n$ sufficiently large,
contrary to our assumption.

For the moment fix $b$. Set $\bar
t_{(a,b)}=t_{(a,b)}'-R_{G_{(a,b)}}(x_{(a,b)})^{-1}$. Since, by supposition the
embedding $\psi_{U_{(a,b)}}$ does not extend  backwards past $\bar t_{(a,b)}$, it
must be the case that $\bar t_{(a,b)}$  is a surgery time and
furthermore that there is a surgery cap ${\mathcal C}_{(a,b)}$ at
this time with the property that there is a point $y_{(a,b)}\in
U_{(a,b)}$ such that $\psi_{U_{(a,b)}}(y_{(a,b)},t)$ converges to a
point $z_{(a,b)}\in {\mathcal C}_{(a,b)}$ as $t$ tends to $\bar
t_{(a,b)}$ from above. () We denote by
$p_{(a,b)}$ the tip of ${\mathcal C}_{(a,b)}$, and we denote by
$\bar h_{(a,b)}$ the scale of the surgery at time $\bar t_{(a,b)}$.



Since the statement that $x_{(a,b)}$ is contained in the core of a $(C,\epsilon)$-cap is
a scale invariant statement, we are free to replace $({\mathcal M}_{(a,b)},G_{(a,b)})$ with
$(\widetilde{\mathcal M}_{(a,b)},\widetilde G_{(a,b)})$, which has been rescaled to make
$\bar h_{(a,b)}=1$ and shifted in time so that $\bar t_{(a,b)}=0$. We denote the new time function
by ${\bf \tilde t}$.
(Notice that this rescaling and time-shifting is different from what we usually do.
Normally, when we have a base point like $x_{(a,b)}$ we rescale to make its scalar curvature one and
we shift time to make it be at time $0$. Here we have rescaled based on the scale of the surgery cap
rather than $R(x_{(a,b)})$.) We set $\tilde Q_{(a,b)}=R_{\widetilde G_{(a,b)}}(x_{(a,b)})$ and we set
$\tilde t'_{(a,b)}={\bf \tilde t}(x_{(a,b)})$. Since the initial time of the strong $\epsilon$-neck
is zero,  $\tilde t'_{(a,b)}=\tilde Q_{(a,b)}^{-1}$. We denote the flow line backward in time from
 $y_{(a,b)}$ by $y_{(a,b)}(\tilde t),\ 0\le \tilde t\le \tilde t'_{(a,b)}$, so that
 $y_{(a,b)}(\tilde t'_{(a,b)})=y_{(a,b)}$.
 Since $U_{(a,b)}$ is a strong $\epsilon$-neck, by our choice of $\epsilon$,
 it follows from Lemma~\ref{neckcurv} and rescaling that
 $R(\psi(y_{(a,b)},\tilde t))$ is within $(0.01) \tilde Q_{(a,b)}$ of
 $\tilde Q_{(a,b)}/(1+\tilde Q_{(a,b)}( \tilde t_{(a,b)}'-\tilde t))$ for all
 $t\in (0,\tilde t_{(a,b)}']$. By taking limits as $t$
 approaches $0$, we see that
  $R_{\widetilde G_{(a,b)}}(z_{(a,b)})$ is within $(0.01)\tilde Q_{(a,b)}$ of $\tilde Q_{(a,b)}/2$.
Let $D$ be the universal constant given in  Lemma~\ref{A_0}, so that
the scalar curvature at any point of the standard initial metric is
at least $D^{-1}$ and at most $D$.  It follows from the third item
in Theorem~\ref{LOCALSURGERY} that, since we have rescaled to make
the surgery scale one, for all $b$ sufficiently large the scalar
curvature on the surgery ${\mathcal C}_{(a,b)}$ is at least
$(2D)^{-1}$ and at most $2D$. In particular, for all $b$
sufficiently large
$$(2D)^{-1}\le R_{\widetilde G_{(a,b)}}(z_{(a,b)})\le 2D.$$
Together with the above estimate relating $R_{\widetilde G_{(a,b)}}(z_{(a,b)})$ and
$\tilde Q_{(a,b)}$, this gives
\begin{equation}\label{DQ}
(5D)^{-1}\le
\tilde Q_{(a,b)}\le 5D.
\end{equation} Since the flow line from
$z_{(a,b)}$ to $y_{(a,b)}$ lies in the closure of a strong
$\epsilon$-neck of scale $\tilde Q_{(a,b)}^{-1/2}$, the scalar curvature is
less than
 $6D$ at every point of this flow line.
 According to Proposition~\ref{Restim} there is $\theta_1<1$
(depending only on $D$) such that $R(q,t)\ge 8D$ for all $(q,t)$
in the standard solution with $t\ge \theta_1$.
% Also, by Lemma~\ref{capball},
% the point $z_{(a,b)}$ is contained in the ball centered at the tip
%of the surgery cap of radius $(A_0+5)\bar h_{(a,b)}$. This implies
%that the limit at time $\bar t_{(a,b)}$  of entire evolving
%$\epsilon$-neck is contained in the ball centered at the tip $p_{(a,b)}$ of
%the cap of radius $(A_0+5)\bar
%h_{(a,b)}+(1.1)\epsilon^{-1}Q_{(a,b)}^{-1/2})$. The above
%inequalities on $Q_{(a,b)}$ give
%$$(A_0+5)\bar h_{(a,b)}+(1.1)\epsilon^{-1}Q_{(a,b)}^{-1/2}
%\le \left(A_0+5+3\epsilon^{-1}\sqrt{D}\right)\bar h_{(a,b)}).$$






By the fifth property of Theorem~\ref{stdsoln} there is
$A'(\theta_1)<\infty$ such that in the
 standard flow, $B(p_0,0,A)$ contains $B(p_0,\theta_1,A/2)$ for every $A\ge A'(\theta_1)$.
 We set $A$ equal to the maximum of $A'(\theta_1)$ and
 $$3\left((1.2)\sqrt{5D}\epsilon^{-1}+(1.1)(A_0+5)+C\sqrt{5D}\right).$$
Now for any $\overline\delta>0$ for all  $b$ sufficiently large, we
have $\delta_{(a,b)}\le \delta''(A,\theta_1,\overline\delta)$, where
$\delta''(A,\theta_1,\overline\delta)$
is the constant given in Proposition~\ref{bigell}.



\begin{claim}[The rescaled surgery time lies before theta_1 in the limit]
\label{claim:tilde-t-prime-bounded-by-theta1}
\uses{def:overline-delta-0}
 Suppose that $b$ is sufficiently large so that
$\delta_{(a,b)}\le \delta''(A,\theta_1,\overline\delta_0)$, where $\overline\delta_0$ is the constant
given in Definition~\ref{defnoverlinedelta0}. Then
$\tilde t_{(a,b)}'\le \theta_1$.
\end{claim}


\begin{proof}
\uses{prop:surgery-cap-evolution-alternative,def:overline-delta-0}
In this proof we shall fix $(a,b)$, so we drop these indices from
the notation.
 Consider
$s\le \theta_1$ maximal so that there is an embedding
$$\psi=\psi_{(a,b)}\colon
B(p_0,0,A)\times [0,s)\to \widetilde{\mathcal
M}_{(a,b)}$$ compatible with time and the vector field. First
suppose that $s<\theta_1$. Then according to
Proposition~\ref{altern} either the entire ball $B(p,0,A)$
disappears at time $s$ or $s$ is the final time of the
time interval of definition for the flow $(\widetilde{\mathcal
M}_{(a,b)},\widetilde G_{(a,b)})$. Since we have the flow line from $z\in
B(p_0,0,A)$ extending to time $\tilde t'=\tilde t'_{(a,b)}$, in either case this
implies that $\tilde t'<s$, proving that $t'<\theta_1$ in this case.

Now suppose that $s=\theta_1$. By the choice of
$\theta_1$, for the standard solution  the scalar curvature at every
$(q,\theta_1)$ is at least $8D$. Since
$\delta_{(a,b)}\le \delta''(A,\theta_1,\overline\delta_0)$, by the
definition of $\overline\delta_0$ given in
Definition~\ref{defnoverlinedelta0} and by Proposition~\ref{altern}
the scalar curvature of the pullback of the metric under  $\psi$ is
within a factor of two of the scalar curvature of the rescaled
standard solution. Hence, the scalar curvature along the flow line
$(z,t)$ through $z$ limits to at least
$8D$ as $t$ tends to $\theta_1$.
Since the scalar curvature on $(z,t)$ for $t\in [0,\tilde t']$ is
bounded above by $6D$, it follows that $\tilde t'<\theta_1$ in this case as
well. This completes the proof of the
claim.
\end{proof}


Thus, we have maps
$$\psi_{(a,b)}\colon B(p_0,0,A)\times [0,\tilde t'_{(a,b)}]\to \widetilde {\mathcal M}_{(a,b)}$$
compatible with time and the vector field, with the property that for each $\delta>0$, for
all $b$ sufficiently large the pullback under this map of $\widetilde G_{(a,b)}$ is within $\delta$
in the $C^{[1/\delta]}$-topology of the restriction of the standard solution.
Let $w_{(a,b)}$ be the result of flowing $x_{(a,b)}$ backward to time $0$.

\begin{claim}[The backward-flowed base point lies in the rescaled standard-flow ball]
\label{claim:w-ab-in-ball-image}
For all $b$ sufficiently large, $w_{(a,b)}\in \psi_{(a,b)}(B(p_0,0,A)\times\{0\})$.
\end{claim}

\begin{proof}
First notice that, by our choice of $\epsilon$,
 every point in the $0$ time-slice of the closure of the strong $\epsilon$-neck
centered at $x_{(a,b)}$ is within distance $(1.1)\tilde Q_{(a,b)}^{-1}\epsilon^{-1}$
of $w_{(a,b)}$. In particular,
$$d_{\widetilde G_{(a,b)}}(w_{(a,b)},y_{(a,b)})<(1.1)\tilde Q_{(a,b)}^{-1/2}\epsilon^{-1}.$$
Since $y_{(a,b)}$ is contained in the surgery cap and the scale of the surgery
at this time is $1$, $y_{(a,b)}$ is within distance $A_0+5$ of $p_{(a,b)}$.
Hence, by the triangle inequality and Inequality~(\ref{DQ}), we have
\begin{eqnarray*}
d_{\widetilde G_{(a,b)}}(w_{(a,b)},p_{(a,b)}) & < & (1.1)\tilde Q_{(a,b)}^{-1/2}\epsilon^{-1}+(A_0+5) \\
& < & (1.1)\sqrt{5D}\epsilon^{-1}+(A_0+5).
\end{eqnarray*}
For $b$ sufficiently large, the image  $\psi_{(a,b)}(B(p_0,0,A))$
contains the ball of radius $(0.95)A$
centered at $p_{(a,b)}$. Since by our choice of $A$ we have $(0.95)A>(1.1)\sqrt{5D}\epsilon^{-1}+(A_0+5)$,
the claim follows.
\end{proof}

We define $q_{(a,b)}\in B(p_0,0,A)$ so that
$\psi_{(a,b)}(q_{(a,b)},0)=w_{(a,b)}$. Of course,
$$\psi_{(a,b)}(q_{(a,b)},\tilde t'_{(a,b)})=x_{(a,b)}.$$ If follows from the
above computation that for all $b$ sufficiently large we have
$$d_0(q_{(a,b)},p_0)<(1.15)\tilde Q_{(a,b)}^{-1/2}\epsilon^{-1}+(1.05)(A_0+5).$$
Since the standard flow has non-negative curvature, it is a distance non-increasing flow. Therefore,
$$d_{\tilde t'_{(a,b)}}(q_{(a,b)},p_0)<(1.15)\tilde Q_{(a,b)}^{-1/2}\epsilon^{-1}+(1.05)(A_0+5).$$

Suppose that a point $(q,\tilde t'_{(a,b)})$ in the standard solution were the center of a
$\beta\epsilon/3$-neck, where $\beta$ is the constant from Proposition~\ref{neckglue}.
 Of course, for all $b$ sufficiently large, $R(q,\tilde
t'_{(a,b)})>(0.99)\tilde Q_{(a,b)}$. Since $\beta<1/2$ and
$\epsilon<\sqrt{5D}(A_0+5)/2$ and $\tilde Q_{(a,b)}\le 5D$, it follows from the
above distance estimate that this neck would contain $(p_0,\tilde t'_{(a,b)})$.
But this is impossible: since $(p_0,\tilde t'_{(a,b)})$ is an isolated fixed
point of an isometric $SO(3)$-action on the standard flow, all the sectional
curvatures at $(p_0,\tilde t'_{(a,b)})$ are equal, and this is in contradiction
with estimates on the sectional curvatures at any point of an $\epsilon$-neck
given in Lemma~\ref{neckcurv}. We can then conclude from Theorem~\ref{stdsolncannbhd}
that for all $b$ sufficiently large, the point $(p_0,\tilde t'_{(a,b)})$ is
contained in the core of a $(C(\beta\epsilon/3),\beta\epsilon/3)$-cap
$Y_{(a,b)}$ in the $\tilde t'_{(a,b)}$ time-slice of the standard solution.
Now note that for
all $b$ sufficiently large, the scalar curvature of $(q_{(a,b)},\tilde
t'_{(a,b)})$ is at least $(0.99)\tilde Q_{(a,b)}$, since the scalar curvature of $x_{(a,b)}$
is equal to $Q_{(a,b)}$. This implies that the
diameter of $Y_{(a,b)}$ is at most
$$(1.01)\tilde Q^{-1/2}_{(a,b)}C(\beta\epsilon/3)<(1.1)\sqrt{5D}C(\beta\epsilon/3).$$
Since $B(p_0,0,A)$ contains $B(p_0,\tilde t'_{(a,b)},A/2)$, and since
$C>C(\beta\epsilon/3)$, it follows from the definition of $A$, the above
distance estimate, and the triangle inequality that for all $b$ sufficiently
large $B(p_0,0,A)\times\{\tilde t'_{(a,b)}\}$ contains $Y_{(a,b)}$.




 Since $C>
C(\beta\epsilon/3)+1$ and since for $b$ sufficiently large
$\psi_{(a,b)}^*\widetilde G_{(a,b)}$ is arbitrarily close to the
restriction of the standard solution metric, it follows from
Lemma~\ref{canonvary} that for all $b$ sufficiently large, the image $\psi_{(a,b)}(Y_{(a,b)})$
 is a $(C,\epsilon)$-cap whose core contains $x_{(a,b)}$.
As we have already remarked, this contradicts the assumption that no $x_n$
has a strong $(C,\epsilon)$-canonical neighborhood.

 This completes the proof in the last case and establishes
Lemma~\ref{16.2}.
\end{proof}






\begin{remark}[The limiting canonical neighborhood is a cap, not a strong epsilon-neck]
\label{rem:limit-canonical-neighborhood-is-cap}
\uses{lem:limit-canonical-neighborhood-fails-at-surgery-scale}
Notice that even though $x_{(a,b)}$ is the center of a strong
$\epsilon$-neck, the canonical neighborhoods of the $x_n$ constructed
in the second case are not a strong $\epsilon$-necks but rather are
$(C,\epsilon)$-caps coming from applying the flow to a
neighborhood of the surgery cap ${\mathcal C}$.
\end{remark}

Now we return to the proof of Proposition~\ref{extend}. For each $a$, we pass
to a subsequence (in $b$) so that Lemma~\ref{16.2} holds for all $(a,b)$. For
each $(a,b)$, let $t_{(a,b)}$ be as in that lemma.  We fix a point
$x_{(a,b)}\in {\bf t}^{-1}(t_{(a,b)})\subset {\mathcal M}_{(a,b)}$ at which the
canonical neighborhood assumption with parameter $r_a$ fails.  For each $a$
choose $b(a)$ such that $\delta_{b(a)}\rightarrow 0$ as $a\rightarrow \infty$. For each
$a$ we set $({\mathcal M}_a,G_a)=({\mathcal M}_{(a,b(a))},G_{(a,b(a))})$, we
set $t_a=t_{(a,b(a))}$, and we let $x_a=x_{(a,b(a))}\in {\mathcal M}_a$. Let
$(\widetilde {\mathcal M}_a,\widetilde G_a)$ be the Ricci flow with surgery
obtained from $({\mathcal M}_a,G_a)$ by shifting $t_a$ to $0$ and rescaling the
metric and time by $R(x_a)$. We have the points $\tilde x_a$ in the $0$
time-slice of $\widetilde {\mathcal M}_a$ corresponding to $x_a\in {\mathcal
M}_a$. Of course, by construction $R_{\widetilde G_a}(\tilde x_a)=1$ for all
$a$.

We shall take limits of a subsequence of this sequence of based
Ricci flows with surgery. Since $r_a\rightarrow 0$ and $R(x_a)\ge
r_a^{-2}$, it follows that $R(x_a)\rightarrow \infty$. By
Proposition~\ref{KAPPALIMIT}, since $\delta_{b(a)}\le \delta(r_a)$
it follows that the restriction of $(\widetilde {\mathcal
M}_a,\widetilde G_a)$ to ${\bf t}^{-1}(-\infty,0)$ is
$\kappa$-non-collapsed on scales $\le \epsilon R_{G_a}^{1/2}(x_a)$.
By passing to a subsequence we arrange that one of the following two
possibilities holds:
\begin{enumerate}
\item[(i)] There is $A<\infty$ and $t'<\infty$ such that,
for each $a$ there is a flow line through a point $y_a$ of $B_{\widetilde
G_a}(\tilde x_a,0,A)$ that is not defined on all of $[-t',0]$. ()
\item[(ii)] For every $A<\infty$ and every $t'<\infty$, for all $a$
sufficiently large all flow lines through points of $B_{\widetilde
G_a}(\tilde x_a,0,A)$ are defined on the interval $[-t', 0]$.
\end{enumerate}




Let us consider the second case. By Proposition~\ref{KAPPALIMIT}
these rescaled solutions are $\kappa$-non collapsed on scales $\le
\epsilon R_{G_a}(x_a)^{1/2}$ for all $t<0$. Since this condition is
a closed constraint, the same is true if $t=0$. Since $R(x_a)\ge
r_a^{-2}$, by construction every point $\tilde
x\in(\widetilde{\mathcal M}_a,\widetilde G_a)$ with $R(\tilde x)\ge
1$ and ${\bf t}(\tilde x)<0$ has a strong $(C,\epsilon)$-canonical
neighborhood.

\begin{claim}[Limit points of scalar curvature exceeding one have a doubled canonical neighborhood]
\label{claim:limit-points-have-canonical-neighborhood} For all $a$ sufficiently large,
every point $\tilde x\in(\widetilde{\mathcal M}_a,\widetilde G_a)$ with $R(\tilde x)>1$ and
${\bf t}(\tilde x)=0$ has a $(2C,2\epsilon)$-canonical neighborhood.
\end{claim}


\begin{proof}
\uses{thm:local-surgery,omegacanon,epslimit}
Assume that $\tilde x\in \widetilde{\mathcal M}$ has $R(\tilde
x)>1$. Suppose that $\tilde x$ is an exposed point. If $a$ is
sufficiently large, then $\delta_{b(a)}$ is arbitrarily close to
zero and hence by the last item in Theorem~\ref{LOCALSURGERY} and
the structure of the standard initial condition,  we see that
$\tilde x$ is contained in the core of a $(2C,2\epsilon)$-cap.

Suppose now that $\tilde x$ is not an exposed point.
 Then we can take a sequence of points
$\tilde y_n\in \widetilde{\mathcal M}_a$ all lying on the flow line for the vector field through $\tilde x$
converging to $\tilde x$ with ${\bf t}(\tilde y_n)<0$. Of course, for all $n$ sufficiently large
$R(\tilde y_n)>1$,
which implies that for all $n$ sufficiently large $\tilde y_n$ has a strong $(C,\epsilon)$-canonical
neighborhood. Passing to a subsequence, we can arrange that all of these canonical neighborhoods
are of the same type. If they are all $\epsilon$-round components, all $C$-components, or all
$(C,\epsilon)$-caps whose cores contain $y_n$, then by taking limits and arguing as in the proof of
Lemma~\ref{omegacanon} we see that $\tilde x$ has a strong
$(2C,2\epsilon)$-canonical neighborhood of the same type. On the other hand, if $\tilde y_n$ is the center
of a strong $\epsilon$-neck for all $n$, then according to Claim~\ref{epslimit},
the limit point $\tilde x$ is the center
of a strong $2\epsilon$-neck.
\end{proof}

Since we have chosen $\epsilon>0$ sufficiently  small so that Theorem~\ref{kaplimit}
applies with $\epsilon$ replaced by $2\epsilon$,  applying this theorem shows
that we can pass to a subsequence and take a smooth limiting
flow of a subsequence of the rescaled flows $(\widetilde {\mathcal
M}_a,\widetilde G_a)$ based at $\tilde x_a$ and defined for all
$t\in (-\infty,0]$. Because the $({\mathcal M}_a,G_a)$ all have
curvature pinched toward positive and since $R(x_a)\rightarrow \infty$
as $a$ tends to infinity, this result says that the limiting flow
has non-negative, bounded curvature and is $\kappa$-non-collapsed on
all scales. That is to say, the limiting flow is a
$\kappa$-solution. By Corollary~\ref{limitcannbhd} this contradicts
the fact that the strong $(C,\epsilon)$-canonical neighborhood
assumption fails at $x_a$ for every $a$. This contradiction shows
that in the second case there is a subsequence of the $a$ such that
$x_a$ has a strong canonical neighborhood and completes the proof of the second case.



Let us consider the first case. In this case we will arrive at a
contradiction by showing that for all $a$ sufficiently large, the
point $x_a$ lies in a strong $(C,\epsilon)$-canonical neighborhood
coming from a surgery cap. Here is the basic result we use to find
that canonical neighborhood.

\begin{lemma}[A bounded-curvature flow line into a surgery cap forces a canonical neighborhood]
\label{lem:std-cap-canonical-neighborhood}\label{lemstdcan}
Suppose that there are $A',D',t'<\infty$ such that the following
holds for all $a$ sufficiently large.  There is a point $y_a\in
B_{\widetilde G_a}(\tilde x_a,0,A')$ and a  flow line of $\chi$
 beginning at $y_a$, defined for backward time and ending at a
 point $z_a$ in a
surgery cap ${\mathcal C}_a$ at time $-t_a$ for some $t_a\le t'$. We
denote this flow line by $y_a(t), -t_a\le t\le 0$. Furthermore,
suppose that the scalar curvature on  the  flow line from $y_a$ to
$z_a$ is bounded by $D'$. Then for all $a$ sufficiently large, $x_a$
has a strong $(C,\epsilon)$-canonical neighborhood.
\end{lemma}


\begin{proof}
\uses{A_0,stdsoln,prop:surgery-cap-evolution-alternative,posdistdec,stdsolncannbhd,canonvary,neckglue}
  The proof is by contradiction. Suppose
the result does not hold. Then there are $A',D',t'<\infty $ and we
can pass to a subsequence (in $a$) such that  the hypotheses of the
lemma hold for every $a$ but no $x_a$ has a strong
$(C,\epsilon)$-canonical neighborhood. The essential point of the
argument is to show that {\bf in the units of the surgery scale} the
elapsed time between the surgery time and $0$ is less than $1$ and
the distance from the point $z_a$ to the tip of the surgery cap is
bounded independent of $a$.

 By Lemma~\ref{A_0}, the fact
that the scalar curvature at $z_a$ is bounded by $D'$ implies that
for all $a$ sufficiently large the scale $\bar h_a$ of the surgery
at time $-t_a$ satisfies
\begin{equation}\label{hequat}
\bar h_a^2\ge (2D'D)^{-1}.
\end{equation}
(Recall that we are working in the rescaled flow $(\widetilde {\mathcal M}_a,\widetilde G_a)$.)



Now we are ready to show that the elapsed time is bounded less than
one in the surgery scale.

\begin{claim}[The elapsed time to the surgery cap is bounded by theta_1 times the surgery scale squared]
\label{claim:elapsed-time-bounded-by-theta1}
There is $\theta_1<1$, depending on $D'$ and $t'$, such that for all
$a$ sufficiently large we have  $t_a< \theta_1\bar h_a^2$.
\end{claim}

\begin{proof}
\uses{prop:surgery-cap-evolution-alternative,def:overline-delta-0,Restim}
We consider two cases: either $t_a\le \bar h_a^2/2$ or $\bar
h_a^2/2<t_a$. In the first case, the claim is obviously true with
$\theta_1$ anything greater than $1/2$ and less than one. In the
second case, the curvature everywhere along the flow line is at most
$ D'<(2t_aD')\bar h_a^{-2}\le (2t'D')\bar h_a^{-2}$. Using
Proposition~\ref{Restim} fix $1/2<\theta_1<1$ so that every point of
the standard solution $(x,t)$ with $t\ge (2\theta_1-1)$ satisfies
$R(x,t)\ge 6t'D'$. Notice that  $\theta_1$ depends only on $D'$ and
$t'$. If $t_a<\theta_1\bar h_a^2$, then the claim holds for this
value of $\theta_1<1$. Suppose $t_a\ge \theta_1\bar h_a^2$, so
that $-t_a+(2\theta_1-1)\bar h_a^2<0$. For all  $a$
sufficiently large we have $\delta_a\le \delta_0''(A_0+5,
\theta_1,\overline \delta_0)$ where $\overline\delta_0$ is the constant
from Definition~\ref{defnoverlinedelta0} and $\delta''_0$ is the constant
from Proposition~\ref{altern}.
This means that the scalar curvatures
at corresponding points of the rescaled standard solution and the
evolution of the surgery cap (up to time $0$) in $\widetilde{\mathcal M}_a$
differ by at most a factor of two. Thus, for these $a$, we have
$R(y_a,(-t_a+(2\theta_1-1)\bar h_a^{2}))\ge 3(t'D')\bar h_a^{-2}$
from the definition of $\overline\delta_0$ and
Proposition~\ref{altern}. But this is impossible since
$-t_a(2\theta_1-1)\bar h_a^3<0$ and $3t'D'/\bar h_a^2\ge 3t_aD'/\bar h_a^2>D'$ as
 $t_a\ge \bar h_a^2/2$. Hence,
$R(y_a,(-t_a+(2\theta_1-1)\bar h_a^2))\le 2t_aD'\bar
h_a^{-2}\le 2t'D'\bar h_a^{-2}$. This contradiction
shows that if $a$ is sufficiently large then $t_a<\theta_1\bar
h_a^2$.
\end{proof}

We pass to a subsequence so that $t_a\bar h_a^{-2}$ converges to
some $\theta\le \theta_1$. We define $\widetilde C$ to be the
maximum of $C$ and $3\epsilon^{-1}\beta^{-1}$. Now, using Part 5 of
Theorem~\ref{stdsoln} we set $A''\ge (9\widetilde
C+3A')\sqrt{2DD'}+6(A_0+5)$ sufficiently large so that in the
 standard flow $B(p_0,0,A'')$ contains
$B(p_0,t,A''/2)$ for any $t\le (\theta_1+1)/2$. This constant is chosen only to
depend on $\theta_1$, $A'$, and $C$. As $a$ tends to infinity, $\delta_a$ tends
to zero which means, by Proposition~\ref{altern}, that for all $a$ sufficiently
large there is an embedding $\rho_a\colon B(p_0,-t_a,A''\bar h_a)\times
[-t_a,0]\to \widetilde {\mathcal M}_a$ compatible with time and the vector
field such that (after translating by $t_a$ to make the flow start at time $0$
and scaling by $\bar h_a^{-2}$) the restriction of $\widetilde G_a$ to this
image is close in the $C^\infty$-topology to the restriction of the standard
flow to $B(p_0,0,A'')\times [0,\bar h_a^{-2}t_a]$. The image $\rho_a(p_0,-t_a)$
is the tip $p_a$ of the surgery cap ${\mathcal C}_a$ in $\widetilde {\mathcal
M}_a$. In particular, for all $a$ sufficiently large the image
$\rho_a\left(B(p_0,-t_a,A''\bar h_a)\times\{0\}\right)$ contains the $A''\bar
h_a/3$-neighborhood of the image $\rho_a(p_0,0)$ of the tip of the surgery cap
under the flow forward to  time $0$. By our choice of $A''$, and
Equation~(\ref{hequat}), this means that for all $a$ sufficiently large
$\rho_a\left(B(p_0,-t_a,A''\bar h_a)\times\{-t_a\}\right)$ contains the
$(3\widetilde C+A') +2(A_0+5)\bar h_a$ neighborhood of $p_a=\rho_a(p_0,-t_a)$.
Notice also that, since the standard solution has positive curvature and hence
the distance between points is non-increasing in time by
Lemma~\ref{posdistdec}, the distance at time $0$ between $\rho_a(p_0,0)$ and
$y_a$ is less than $2(A_0+5)\bar h_a$. By the triangle inequality, we conclude
that for all $a$ sufficiently large, $\rho_a\left(B(p_0,-t_a,A''\bar h_a)\times
\{0\}\right)$ contains the $3\widetilde C$-neighborhood of $x_a$. Since the
family of metrics on $\rho_a\left(B(p_0,-t_a,A''\bar h_a)\times
[-t_a,0]\right)$ (after time-shifting by $t_a$ and rescaling by $\bar
h_a^{-2}$) are converging smoothly to the ball $B(p_0,0,A'')\times [0,\theta]$
in the standard flow,
 for all $a$ sufficiently large then the flow from time $-t_a$ to $0$ on the
$3\widetilde C$-neighborhood of $x_a$ is, after rescaling by $\bar
h_a^{-2}$,
 very nearly isometric to the restriction of the standard flow from time $0$ to $\bar h^{-2}_at_a$ on
 the $3\widetilde C\bar h_a^{-1}$-neighborhood of some
point $q_a$ in the standard flow. Of course, since the scalar
curvature of $x_a$ is one,  $R(q_a,\bar h_a^{-2}t_a)$ in the standard flow is
close to $\bar h_a^{-2}$. Hence, by Theorem~\ref{stdsolncannbhd}
there is a neighborhood $X$ of $(q_a,\bar h_a^{-2}t_a)$ in the standard solution
that either is a $(C,\epsilon)$-cap, or is an evolving
$\beta\epsilon/3$-neck centered at $(q_a,\bar h_a^{-2}t_a)$. In the latter case
either the evolving neck is defined for backward time
$(1+\beta\epsilon/3)$ or its initial time-slice is the zero
time-slice and this initial time-slice lies at distance at least $1$
from the surgery cap. Of course, $X$ is contained in the
$CR(q_a,\bar h_a^{-2}t_a)^{-1/2}$ neighborhood of $(q_a,\bar h_a^{-2}t_a)$ in the standard
solution. Since $\widetilde C\ge C$ and $R(q_a,\bar h_a^{-2}t_a)$ is close to
$\bar h_a^{-2}$, the neighborhood $X$ is contained in the
$2\widetilde C\bar h_a^{-1}$-neighborhood of $(q_a,\bar h_a^{-2}t_a)$ in the
standard solution.  Hence, after rescaling, the corresponding
neighborhood of $x_a$ is contained in the $3\widetilde
C$-neighborhood of $x_a$. If either of the first two cases in Theorem~\ref{stdsolncannbhd} occurs
for a subsequence of $a$ tending to infinity, then  by
Lemma~\ref{canonvary} and the fact that $\widetilde C>\mathrm{max}(C,\epsilon^{-1})$, we see that there is a subsequence of $a$
for which $x_a$ either is contained in the core of a
$(C,\epsilon)$-cap or is the center of a strong $\epsilon$-neck.

We must examine further the last case. We suppose that for every $a$
this last case holds. Then for all $a$ sufficiently large we have an
$\beta\epsilon/3$-neck $N_a$ in the zero time-slice of $\widetilde{\mathcal
M}_a$ centered at $x_a$. It is an evolving neck and there is an
embedding $\psi\colon N_a\times [-t_a,0]\to \widetilde{\mathcal
M}_a$ compatible with time and the vector field so that the initial
time-slice $\psi(N_a\times \{-t_a\})$ is in the surgery time-slice
$M_{-t_a}$ and is disjoint from the surgery cap, so in fact it is
contained in the continuing region at time $-t_a$. As we saw above,
the image of the central $2$-sphere $\psi(S^2_a\times\{-t_a\})$ lies
at distance at most  $ A''\bar h_a$ from the tip of the surgery cap
$p_a$ (where, recall, $A''$ is a constant independent of $a$). The
$2$-sphere, $\Sigma_a$, along which we do surgery,   creating the surgery
cap with $p_a$ as its tip, is the central $2$-sphere of a strong
$\delta_{b(a)}$-neck. As $a$ tends to infinity the surgery control
parameter $\delta_{b(a)}$ tends to zero. Thus, for $a$ sufficiently
large this strong $\delta_{b(a)}$-neck will contain a strong
$\beta\epsilon/2$- neck $N'$ centered at $\psi(x_a,-t_a)$. Since we
know that the continuing region at time $-t_a$ contains a
$\beta\epsilon/3$-neck centered at $(x_a,-t_a)$, it follows that
$N'$ is also contained in $C_{-t_a}$. That is to say, $N'$  is
contained in the negative half of the $\delta_{b(a)}$-neck centered
at $\Sigma_a$. Now we are in the situation of
Proposition~\ref{neckglue}. Applying this result tells us that $x_a$
is the center of a strong $\epsilon$-neck.

This completes the proof that for all $a$ sufficiently large, $x_a$
has a $(C,\epsilon)$-canonical neighborhood in contradition to our
assumption. This contradiction completes the proof of
Lemma~\ref{lemstdcan}.
\end{proof}




There are several steps required to complete the proof of
Proposition~\ref{extend}. The first step helps us apply the previous claim to
find strong $(C,\epsilon)$-canonical neighborhoods.


\begin{claim}[Uniform curvature bound along backward flow lines from a fixed ball]
\label{claim:curvature-bound-along-backward-flow-lines}\label{D(A)delta(A)}
Given any $A<\infty$ there is $D(A)<\infty$ and $\delta(A)>0$ such that for all $a$
sufficiently large, $|\Rm|$ is bounded by $D(A)$ along all backward flow lines
beginning at a point of $B_{\widetilde
G_a}(\tilde x_a,0,A)$ and defined for backward time at most $\delta(A)$.
\end{claim}

\begin{proof}
\uses{prop:kappa-limit-noncollapsing,bcbd,controlnbhd}
Since all points  $y\in ({\mathcal M}_a,G_a)$ with $R_{G_a}(y)\ge
r^{-2}_a$ and ${\bf t}(y)<{\bf t}(x_a)$ have strong
$(C,\epsilon)$-canonical neighborhoods, and since $R(x_a)=r^{-2}_a$,
we see that all points $y\in (\widetilde{\mathcal M}_a,\widetilde
G_a)$ with ${\bf t}(y_a)<0$ and with $R_{\tilde G_a}(y_a)\ge 1$ have
strong $(C,\epsilon)$-canonical neighborhoods. It follows that all
points in $(\widetilde {\mathcal M}_a,\widetilde G_a)$ with ${\bf
t}(y)\le 0$ and $R(y)>1$ have strong $(2C,2\epsilon)$-canonical
neighborhoods. Also, since $\delta_a\le \delta(r_a)$, where
$\delta(r_a)$ is the constant given in Proposition~\ref{KAPPALIMIT},
and since the condition of being $\kappa$-non-collapsed is a closed
constraint, it follows from Proposition~\ref{KAPPALIMIT} that these
Ricci flows with surgery are $\kappa$-non-collapsed for a fixed
$\kappa>0$. It is now immediate from Theorem~\ref{bcbd} that there
is a constant $D_0(A)$ such that $R$ is bounded above on
$B_{\widetilde G_a}(\tilde x_a,0,A)$ by $D_0(A)$. Since every point
$y\in ({\mathcal M}_a,G_a)$ with $R(y)>1$ of the sequence of with
scalar curvature at least $1$ has a $(C,\epsilon)$ canonical
neighborhood, it follows from the definition that for every such
point $y$ we have $\left|\partial R(y)/\partial t\right|<CR(y)^2$.
Arguing as in Lemma~\ref{controlnbhd} we see that there is a
constant $\delta(A)>0$ and a bound $D'(A)$, both depending only in
$D_0(A)$,
 for the scalar curvature at all points
of backward flow lines beginning in $B_{\widetilde
G_a}(\tilde x_a,0,A)$ and defined for backward time at most $\delta(A)$. Since the curvature is pinched toward positive,
it follows that there is a bound $D(A)$ depending only on $D'(A)$ to $|\Rm|$ on the same flow lines.
\end{proof}

\begin{claim}[Dichotomy: uniform backward parabolic neighborhood or a canonical neighborhood]
\label{claim:parabolic-neighborhood-or-canonical-neighborhood} After passing to a subsequence (in $a$), either:
\begin{enumerate}
\item[(1)]
for each $A<\infty$ there are $D(A)<\infty$ and $t(A)>0$ such that for all $a$
sufficiently large $P_{\widetilde G_a}(\tilde x_a,0,A,-t(A))$ exists in
$\widetilde {\mathcal M}_a$ and $|\Rm|$ is bounded by $D(A)$ on this
backward parabolic neighborhood, or
\item[(2)] each $x_a$ has a strong
$(C,\epsilon)$-canonical neighborhood.
\end{enumerate}
\end{claim}


\begin{proof}
\uses{claim:curvature-bound-along-backward-flow-lines,lem:std-cap-canonical-neighborhood}
First notice that if there is $t(A)>0$ for which the backwards parabolic neighborhood
$P=P_{\widetilde G_a}(\tilde x_a,0,A,-t(A))$ exists,
then, by Claim~\ref{D(A)delta(A)},
there are constants $D(A)<\infty$ and $\delta(A)>0$ such that, replacing $t(A)$ by $\mathrm{min}(t(A),\delta(A))$,
 $|\Rm|$ is bounded by $D(A)$ on $P$.
Thus, either Item (1) holds or
passing to a subsequence, we can suppose that there is some $A<\infty$ for which no $t(A)>0$ as
required by Item (1) exists.
Then,  for each $a$ we find a point
$y_a\in B_{\widetilde G_a}(\tilde x_a,0,A)$ such that the backwards
flow line from $y_a$ meets a surgery cap at a time $-t_a$ where
$\mathrm{lim}_{a\rightarrow \infty}(t_a)=0$. Then, by the previous claim, for all $a$ sufficiently large, the
sectional curvature along any backward flow line beginning in $B_{\widetilde G_a}(\widetilde x_a,0,A)$
and defined for backward time $t_a$ is bounded by
a constant $D(A)$ independent of $a$.  Under our assumption this
means that for all $a$ sufficiently large, there is a point $y_a\in
B_{\widetilde G_a}(\tilde x_a,0,A)$ and a backwards flow line
starting at $y_a$ ending at a point $z_a$ of a surgery cap, and the
sectional curvature along this entire flow line is bounded by $D(A)<\infty$. Thus,
applying Lemma~\ref{lemstdcan} produces the strong
$(C,\epsilon)$-canonical neighborhood around $x_a$, proving the claim.
\end{proof}


But we are assuming that no $x_a$ has a strong $(C,\epsilon)$-canonical
neighborhood. Thus, the consequence of the previous claim is that for each
$A<\infty$ there is a $t(A)>0$ such that for all $a$ sufficiently large
$P_{\widetilde G_a}(\tilde x_a,0,A,-t(A))$ exists in $\widetilde {\mathcal
M}_a$ and there is a bound, depending only on $A$ for $|\Rm|$ on
this backward parabolic neighborhood. Applying Theorem~\ref{sliceconv} we see that, after passing to a
subsequence, there is a smooth limit $(M_\infty,g_\infty,x_\infty)$ to the zero
time-slices $(\widetilde{\mathcal M}_a,\widetilde G_a,\tilde x_a)$.
Clearly, since the curvatures of the sequence are pinched toward positive, this
limit has non-negative curvature.

Lastly, we show that $(M_\infty,g_\infty)$ has bounded curvature. By Part 3 of Proposition~\ref{canonvary}
each point of $(M_\infty,g_\infty)$ with scalar curvature
greater than one has a $(2C,2\epsilon)$-canonical neighborhood. If a point lies in an $2\epsilon$-component
or in a $2C$-component, then $M_\infty$ is compact, and hence clearly has bounded curvature.
Thus, we can assume that each $y\in M_\infty$ with $R(y)>1$ is either the center of a $2\epsilon$-neck
or is contained in the core of a $(2C,2\epsilon)$-cap. According to Proposition~\ref{narrows}
 $(M_\infty,g_\infty)$ does not contain $2\epsilon$-necks of arbitrarily high curvature.
 It now follows then that $(M_\infty,g_\infty)$ there is a bound to the scalar curvature of any
 $2\epsilon$-neck and of any $(2C,2\epsilon)$-cap, and hence it follows that $(M_\infty,g_\infty)$
 has bounded curvature.


\begin{claim}[If the backward time bound cannot be made uniform in the radius, canonical neighborhoods exist]
\label{claim:uniform-t-prime-independent-of-A}
If the constant $t(A)>0$ cannot be chosen independent of $A$, then
after passing to a subsequence, the $x_a$ have strong
$(C,\epsilon)$-canonical neighborhoods.
\end{claim}

\begin{proof}
\uses{lem:std-cap-canonical-neighborhood,controlnbhd}
 Let $Q$ be the bound of the scalar curvature of
$(M_\infty,g_\infty,x_\infty)$. Then by Lemma~\ref{controlnbhd}
 there is a constant
$\Delta t>0$ such that if $R_{\widetilde G_a}(y,0)\le 2Q$, then the
scalar curvature is bounded by $16Q$ on the backward flow line from
$y$ defined for any time $\le \Delta t$. Suppose that there is
$A<\infty$ and a  subsequence of $a$ for which there is a flow line
beginning at a point $y_a \in B_{\widetilde G_a}(\tilde x_a,0,A)$
defined for backward time at most $\Delta t$ and ending at a point
$z_a$ of a surgery cap. Of course, the fact that the scalar
curvature of $(M_\infty,g_\infty)$ is at most $Q$ implies that for
all $a$ sufficiently large, the scalar curvature of $B_{\widetilde
G_a}(\tilde x_a,0,A)$ is less than $2Q$. This implies  that for all
$a$ sufficiently large the scalar curvature along the flow line
from $y_a$ to $z_a$ in a surgery cap is $\le 16Q$. Now invoking
Lemma~\ref{lemstdcan} we see that for all $a$ sufficiently large the
point $\widetilde x_a$ has a strong $(C,\epsilon)$-canonical
neighborhood. This is a contradiction, and this contradiction proves
that we can choose $t(A)>0$ independent of $A$.
\end{proof}


Since we are assuming that no $x_a$ has a strong
$(C,\epsilon)$-canonical neighborhood, this means that it is
possible to find a constant $t'>0$ such that $t(A)\ge t'$ for all
$A<\infty$. Now let $0<T'\le \infty$ be the maximum possible value for such
$t'$. Then for every $A$ and every $T<T'$ the
parabolic neighborhood $P_{\widetilde G_a}(\tilde x_a,0,A,T)$ exists
for all $a$ sufficiently large. According to
Theorem~\ref{kaplimit}, after passing to a subsequence, there is
a limiting flow $(M_\infty,g_\infty(t),x_\infty),\ -T'<t\le 0$, and
this limiting flow has bounded, non-negative curvature.
If $T=\infty$, this limit is a $\kappa$-solution, and hence the $x_a$ have
strong $(C,\epsilon)$-canonical neighborhoods for all $a$ sufficiently large,
which is a contradiction.

Thus, we can assume that $T'<\infty$.
 Let $Q$ be
the bound for the scalar curvature of this flow. Since $T'$ is
maximal, for every $t>T'$, after passing to a subsequence, for all
$a$ sufficiently large there is $A<\infty$ and a backwards flow
line, defined for a time less than $t$, starting at a point $y_a$ of
$B_{\widetilde G_a}(\tilde x_a,0,A)$ and ending at a point $z_a$ of
a surgery cap. Invoking Lemma~\ref{controlnbhd} again, we see that
for all $a$ sufficiently large, the scalar curvature is bounded on
the flow line from $y_a$ to $z_a$ by a constant independent of $a$.
Hence, as before, we see that for all $a$ sufficiently large $x_a$
has a strong $(C,\epsilon)$-canonical neighborhood; again this is  a contradiction.

Hence, we have now shown  that our assumption that the strong
$(C,\epsilon)$-canonical neighborhood assumption fails for all $r_a$
and all $\delta_{a,b}$ leads to a contradiction and hence is false.


This completes the proof of Proposition~\ref{extend}.
\end{proof}


\section{Surgery times don't accumulate}\label{sectcomplete}


Now we turn to the proof of Theorem~\ref{MAIN}.
 Given surgery parameter
sequences $$\Delta_i=\{\delta_0,\ldots,\delta_i\}$$
$${\bf r_i}=\{r_0,\ldots,r_i\}$$
$${\bf K_i}=\{\kappa_0,\ldots,\kappa_i\},$$
 we let
$r_{i+1}$ and $\delta_{i+1}$ be as in Proposition~\ref{extend} and
then set $\kappa_{i+1}=\kappa(r_{i+1})$ as in
Proposition~\ref{KAPPALIMIT}.
 Set $${\bf
r_{i+1}}=\{{\bf r_i},r_{i+1}\}$$
$${\bf K_{i+1}}=\{{\bf
K_i},\kappa_{i+1}\}$$
$$\Delta_{i+1}=\{\delta_0,\ldots,\delta_{i-1},\delta_{i+1},\delta_{i+1}\}.$$
Of course, these are also surgery parameter sequences.

Let $\overline \delta\colon [0,T]\to \Ar^+$ be any non-increasing positive
function and let $({\mathcal M},G)$ be a Ricci flow with surgery defined on
$[0,T)$ for some $T\in [T_i,T_{i+1})$ with surgery control parameter
$\overline\delta$. Suppose $\overline \delta\le \Delta_{i+1}$ and that this
Ricci flow with surgery satisfies the conclusion of Theorem~\ref{MAIN} with
respect to these sequences on its entire interval of definition. We wish to
extend this Ricci flow with surgery to one defined on $[0,T')$ for some $T'$
with $T<T'\le T_{i+1}$ in such a way that $\overline\delta$ is the surgery
control parameter and the extended Ricci flow with surgery continues to satisfy
the conclusions of Theorem~\ref{MAIN} on its entire interval of definition.

We may as well assume that the Ricci flow $({\mathcal M},G)$ becomes
singular at time $T$. Otherwise we would simply extend by Ricci flow
to a later time $T'$. By Proposition~\ref{KAPPALIMIT} and
Proposition~\ref{extend} this extension will continue to satisfy the
conclusions of Theorem~\ref{MAIN} on its entire interval of
definition. If $T\ge T_{i+1}$, then we have extended the Ricci flow
with surgery to time $T_{i+1}$ as required and hence completed the
inductive step. Thus, we may as well assume that $T<T_{i+1}$.

Consider the maximal extension of $({\mathcal M},G)$ to time $T$. Let $T^-$ be
the previous surgery time, if there is one, and otherwise be zero. If the $T$
time-slice, $\Omega(T)$, of this maximal extension is all of $M_{T^-}$, then
the curvature remains bounded as $t$ approaches $T$ from below. According to
Proposition~\ref{flowextend} this means that $T$ is not a surgery time and we
can extend the Ricci flow on $(M_{T^-},g(t)),\ T^-\le t<T$, to a Ricci flow on
$(M_{T^-},g(t)),\ T^-\le t<T'$ for the maximal time interval (i.e. so that the
flow becomes singular at time $T'$ or $T'=\infty$). But we are assuming that
the flow goes singular at  $T$. That is to say,  $\Omega(T)\not= M_{T^-}$. Then
we can do surgery at time $T$ using $\overline \delta(T)$ as the surgery
control parameter, setting $\rho(T)=r_{i+1}\delta(T)$. Let $(M_T,G(T))$ be the
result of surgery. If $\Omega_{\rho(T)}(T)=\emptyset$, then the surgery process
at time $T$ removes all of $M_{T'}$. In this case, the Ricci flow is understood
to exist for all time and to be empty for $t\ge T$. In this case we have
completed the extension to $T_{i+1}$, and in fact all the way to $T=\infty$,
and hence completed the inductive step in the proof of the proposition.

We may as well assume that $\Omega_{\rho(T)}(T)\not=\emptyset$ so
that the result of surgery is a non-empty manifold $M_T$.
 Then we use this compact Riemannian
$3$-manifold as the initial conditions of a Ricci flow beginning
at time $T$. According to Lemma~\ref{14.26} the union along
$\Omega(T)$ at time $T$ of this Ricci flow with $({\mathcal M},G)$
is a Ricci flow with surgery satisfying Assumptions (1) -- (7) and whose
curvature is pinched toward positive.

Since the surgery  control parameter $\overline\delta(t)$ is at most
$\delta(r_{i+1})$, the constant from Proposition~\ref{KAPPALIMIT},
 for all $t\in [T_{i-1},T]$, since $T\le T_{i+1}$,
and since the restriction of $({\mathcal M},G)$ to ${\bf
t}^{-1}([0,T_i))$ satisfies Proposition~\ref{KAPPALIMIT}, we see by
Proposition~\ref{extend} that the extended Ricci flow with surgery
satisfies the conclusion of Theorem~\ref{MAIN} on its entire time
interval of definition.

Either we can repeatedly apply this process, passing from one surgery
time to the next and eventually reach $T\ge T_{i+1}$, which would
prove the inductive step, or there is an unbounded number of
surgeries in the time interval $[T_i,T_{i+1})$. We must rule out the
latter case.

\begin{lemma}[The number of surgery times up to T_{i+1} is bounded]
\label{lem:bounded-number-of-surgeries}
\uses{thm:main-existence-ricci-flow-surgery}\label{bddsurgery}
Given a Ricci flow with surgery $({\mathcal M},G)$ defined on $[0,T)$ with
$T\le T_{i+1}$ with surgery control parameter $\overline\delta$ a
non-increasing positive function defined on $[0,T_{i+1}]$ satisfying the
hypotheses of Theorem~\ref{MAIN} on its entire time-domain of definition, there
is a constant $N$ depending only on the volume of $(M_0,g(0))$, on $T_{i+1}$,
on $r_{i+1}$, and on $\overline\delta(T_{i+1})$ such that this Ricci flow with
surgery defined on the interval $[0,T)$ has at most $N$ surgery times.
\end{lemma}


\begin{proof}
\uses{A_0}
Let $(M_t,g(t))$  be the $t$ time-slice of $({\mathcal M},G)$. If $t_0$ is not
a surgery time, then $\Vol(t)=\Vol(M_t,g(t))$ is a smooth function
of $t$ near $t_0$ and
$$\frac{d\Vol}{dt}(t_0)=-\int_{M_{t_0}}Rd\vol,$$ so that, because of the curvature
pinching toward positive hypothesis, we have $\frac{d\Vol}{dt}(t_0)\le
6\Vol(t_0)$. If $t_0$ is a surgery time, then either $M_{t_0}$ has fewer
connected components than  $M_{t_0^-}$ or we do a surgery in an $\epsilon$-horn
of $M_{t_0^-}$. In the latter case we remove the end of the $\epsilon$-horn,
which  contains the positive half of a $\overline\delta(t_0)$-neck of scale
$h(t_0)$. We then  sew in a ball with volume at most $(1+\epsilon)Kh^3(t_0)$,
where $K<\infty$ is the universal constant given in Lemma~\ref{A_0}. Since
$h(t_0)\le \overline\delta^2(t_0)r(t_0)\le \delta_0^2r(t_0)$ and since we have
chosen
 $\overline\delta(t_0)\le\delta_0<K^{-1}$, it follows that this
operation lowers volume by at least $\delta^{-1}h^2(t_0)/2$. Since
$\overline\delta(t_0)\ge \overline \delta(T_{i+1})>0$ and the canonical
neighborhood parameter $r$ at time $t_0$ is at least $r_{i+1}>0$, it follows
that $h(t_0)\ge h(T_{i+1})>0$. Thus, each surgery at time $ t_0\le T_{i+1}$
along a $2$-sphere removes at least a fixed amount of volume depending on
$\overline \delta(T_{i+1})$ and $r_{i+1}$. Since under Ricci flow the volume
grows at most exponentially, we see that there is a bound depending only on
$\overline\delta(T_{i+1})$, $T_{i+1}$, $r_{i+1}$ and $\Vol(M_0,g(0))$ to
the number of $2$-sphere surgeries that we can do in this time interval. On the
other hand, the number of components at any time $t$ is at most $N_0+S(t)-D(t)$
where $N_0$ is the number of connected components of $M_0$, $S(t)$ is the
number of $2$-sphere surgeries performed in the time interval $[0,t)$ and
$D(t)$ is the number of connected components removed by surgeries at times in
the interval $[0,t)$.  Hence, there is a bound on the number of components in
terms of $N_0$ and $S(T)$ that can be removed by surgery in the interval
$[0,T)$. Since the initial conditions are normalized, $N_0$ is bounded by the
volume of  $(M_0,g(0))$. This completes the proof of the result.
\end{proof}


This lemma completes the proof of the fact that for any $T\le T_{i+1}$, we
encounter only a fixed bounded number  surgeries in  the Ricci flow with
surgery from $0$ to $T$. The bound depends on the volume of the initial
manifold as well as the surgery constants up to time $T_{i+1}$. In particular,
for a given initial metric $(M_0,g(0))$ there is a uniform bound, depending
only on the surgery constants up to time $T_{i+1}$, on the number of surgeries
in any Ricci flow with surgery defined on a subinterval of $[0,T_{i+1})$. It
follows that the surgery times cannot accumulate in any finite interval. This
completes the proof of Theorem~\ref{MAIN}.


To sum up, we have sequences $\Delta$, ${\bf K}$ and ${\bf r}$ as given in
Theorem~\ref{MAIN}. Let $\overline \delta\colon [0,\infty)\to \Ar$ be a
positive, non-increasing function with $\overline \delta\le \Delta$. Let $M$ be
a compact $3$-manifold that contains no embedded $\Ar P^2$ with trivial normal
bundle. We have proved that for any normalized initial Riemannian metric
$(M_0,g_0)$ there is a Ricci flow with surgery with time-interval of definition
$[0,\infty)$ and with $(M_0,g_0)$ as initial conditions. This Ricci flow with
surgery is ${\bf K}$-non-collapsed and satisfies the strong
$(C,\epsilon)$-canonical neighborhood theorem with respect to the parameter
${\bf r}$. It also has curvature pinched toward positive. Lastly, for any $T\in
[0,\infty)$ if there is a surgery at time $T$ then this surgery is performed
using the surgery parameters $\overline\delta(T)$ and $r(T)$, where if $T\in
[T_i,T_{i+1})$ then $r(T)=r_{i+1}$. In this Ricci flow with surgery, there are
only finitely many surgeries on each finite time interval. As far as we know
there may be infinitely many surgeries in all.
